\begin{exercise}\label{ex:o1oval}
要用一个折射界面把介质$n_1$中的点$A$严格成像到介质$n_2$中的点$B$。忽略反射，证明等光程曲面$n_1|Q-A|+n_2|B-Q|=C$的适当光滑部分满足折射定律，说明与球面近轴成像的区别。
\Diagram{fig:o1ovalq}{笛卡尔卵形面的一个分支，图取$n_1=1,n_2=1.5,C=7.5$}{\fill(0,0)circle(1.5pt)node[below]{$A$};\fill(5.5,0)circle(1.5pt)node[below]{$B$};\draw[curve]plot coordinates {(3.732186,-1.500000) (3.532345,-1.462025) (3.375574,-1.424051) (3.242276,-1.386076) (3.124437,-1.348101) (3.017837,-1.310127) (2.919932,-1.272152) (2.829045,-1.234177) (2.744007,-1.196203) (2.663963,-1.158228) (2.588269,-1.120253) (2.516428,-1.082278) (2.448047,-1.044304) (2.382811,-1.006329) (2.320466,-0.968354) (2.260808,-0.930380) (2.203667,-0.892405) (2.148908,-0.854430) (2.096420,-0.816456) (2.046116,-0.778481) (1.997930,-0.740506) (1.951812,-0.702532) (1.907731,-0.664557) (1.865669,-0.626582) (1.825624,-0.588608) (1.787608,-0.550633) (1.751645,-0.512658) (1.717772,-0.474684) (1.686040,-0.436709) (1.656510,-0.398734) (1.629252,-0.360759) (1.604347,-0.322785) (1.581883,-0.284810) (1.561951,-0.246835) (1.544645,-0.208861) (1.530055,-0.170886) (1.518266,-0.132911) (1.509353,-0.094937) (1.503375,-0.056962) (1.500375,-0.018987) (1.500375,0.018987) (1.503375,0.056962) (1.509353,0.094937) (1.518266,0.132911) (1.530055,0.170886) (1.544645,0.208861) (1.561951,0.246835) (1.581883,0.284810) (1.604347,0.322785) (1.629252,0.360759) (1.656510,0.398734) (1.686040,0.436709) (1.717772,0.474684) (1.751645,0.512658) (1.787608,0.550633) (1.825624,0.588608) (1.865669,0.626582) (1.907731,0.664557) (1.951812,0.702532) (1.997930,0.740506) (2.046116,0.778481) (2.096420,0.816456) (2.148908,0.854430) (2.203667,0.892405) (2.260808,0.930380) (2.320466,0.968354) (2.382811,1.006329) (2.448047,1.044304) (2.516428,1.082278) (2.588269,1.120253) (2.663963,1.158228) (2.744007,1.196203) (2.829045,1.234177) (2.919932,1.272152) (3.017837,1.310127) (3.124437,1.348101) (3.242276,1.386076) (3.375574,1.424051) (3.532345,1.462025) (3.732186,1.500000)};\draw[curve](0,0)--(2.372225,1)--(5.5,0);\draw[guide](0,0)--(2.372225,-1)--(5.5,0);\node[above]at(2.372225,1){$Q$};}
\end{exercise}
\begin{solution}
定义$F(Q)=n_1r_1+n_2r_2$，沿入射传播方向$\vct t_1=(Q-A)/r_1$、出射方向$\vct t_2=(B-Q)/r_2$。梯度为$\nabla F=n_1\vct t_1-n_2\vct t_2$。图\ref{fig:o1ovalq}的等值面法线沿$\nabla F$，任何切向位移$\delta\vct q$满足$\nabla F\cdot\delta\vct q=0$，所以$n_1\vct t_{1\parallel}=n_2\vct t_{2\parallel}$。
\Diagram{fig:o1ovala}{等光程面的法线由入射、出射约化方向矢量之差确定}{\draw[axis](0,0)--(2.4,1.2)node[above left]{$n_1\vct t_1$};\draw[axis](0,0)--(2.1,-.5)node[below]{$n_2\vct t_2$};\draw[->,thick](2.1,-.5)--(2.4,1.2);\node[right]at(2.6,.35){差矢量};\draw[guide](0,0)--(.3,1.7);}
上式为斯涅尔切向条件，法线的矢量构造见图\ref{fig:o1ovala}。须选择使光由介质1进入介质2、无全反射且不被其他界面截断的分支。球面一般只能在轴上小孔径处近似满足固定物像等光程；这个等光程面通常不是球面，可以为指定物像对校正球差，却不保证别的物点也严格成像。
\end{solution}

\begin{exercise}\label{ex:o1galilean}
空气中正物镜焦距$f_1>0$，负目镜$f_2<0$，间隔$d=f_1+f_2>0$。求无焦系统角放大率、轴向平行束宽度比和像的正倒，与双正透镜系统比较。
\Diagram{fig:o1galileoq}{负目镜截取尚未到物镜焦点的会聚光束}{\OptAxis\draw[<->](1,-1.3)--(1,1.3);\draw[>-<](4,-1)--(4,1);\draw[guide](5.5,-1.5)--(5.5,1.5);\node[above]at(5.5,1.5){物镜焦点};\node[above]at(4,1){负目镜};}
\end{exercise}
\begin{solution}
图\ref{fig:o1galileoq}仍用两透镜矩阵，把$d=f_1+f_2$代入得到
\[
 M=\begin{pmatrix}-f_2/f_1&f_1+f_2\\0&-f_1/f_2\end{pmatrix}.
\]
因$f_2<0$，两个对角元都正。轴向平行光的宽度比$|f_2|/f_1$，小视角放大率$\theta_o/\theta_i=f_1/|f_2|>0$，所以角方向不反转。
\Diagram{fig:o1galileoa}{光束在到达几何焦点之前被负透镜重新准直}{\OptAxis\draw[<->](1,-1.3)--(1,1.3);\draw[>-<](4,-1)--(4,1);\draw[->](0,1)--(1,1)--(4,.333333)--(6,.333333);\draw[->](0,-1)--(1,-1)--(4,-.333333)--(6,-.333333);\draw[guide](4,.333333)--(5.5,0);\draw[guide](4,-.333333)--(5.5,0);}
图\ref{fig:o1galileoa}按$f_1=4.5,f_2=-1.5,d=3$作图，束宽缩为1/3、角放大率3，内部没有经过实焦点。双正系统角放大率为负，会产生倒向角像。两系统都满足$\Phi=0$，但角放大率符号和内部光路不同，对应正立、倒立两种观察效果。
\end{solution}
